\scanpage{277}


d'équivalence \deleted{\unsure{entre}}\inserted{des} $\unsure{t|f|S}$

\[
t\sim t'\ \overset{\mathrm{déf}}{\Longleftrightarrow}\ \operatorname{Et}(s)=\operatorname{Et}(t').
\]

Cette dernière propriété est une \deleted{\unsure{con}}caractérisation de $\underline a$ commune à tous les cas (aussi bien $n=4$ que $n\ne4$).

2) Soient $\Lambda$ un anneau semi-local, \deleted{et soient}\inserted{$M$ un} module libre de rang 2 sur $\Lambda$, $u,u'\in M^*$ tels que l'on ait

\[
\left\{\begin{aligned}
&\forall\,x\in M^*,\qquad u\wedge x\in(\det M)^*\ \Longleftrightarrow
\bigl(\deleted{u'\wedge x\in(\det M)^*}\\
&\hspace{14em}\deleted{\Longrightarrow}\inserted{\text{et}}\quad u'\wedge x\in\pm\{u\wedge x\}
\end{aligned}\right.\tag*{$(*)$}
\]

% slanted
\margin{\uline{NB} Le cas exceptionnel est aussi celui où le normalisateur de $\operatorname{Gl}_\Lambda(M)^{\pm}$ dans $\unsure{\mathfrak S_{S_M}|\mathcal G_{S_M}}=\mathfrak S_{(M^*/\pm1)}$ n'est pas réduit au groupe \deleted{\unsure{$Z$}} $\Lambda^*/\pm1$ opérant de façon \unsure{identique|évidente} et par cons, \unsure{l'antipodisme} \unsure{fleurit|fleuris}.}

Alors, on a $u'\in\{\pm u\}$, sauf dans le cas suivant :

\[
\begin{aligned}
&\deleted{\unsure{2_\Lambda\in\Lambda}}\quad 2\Lambda=\{0,2\}\ne\inserted{\{0\}}\deleted{\unsure{\Lambda}}\ \text{i.e.}\\[-2pt]
&\qquad\left\{\begin{aligned}
&\inserted{\Lambda\ \text{une }\mathbb Z/4\mathbb Z\text{-alg.}}\\
&\deleted{\unsure{2_\Lambda\ne0,\ 2_\Lambda\in\operatorname{Nil}(\Lambda)\ \text{et}}}\ \text{extension}\\
&\text{d'une }\mathbb F_2\text{-algèbre }\Lambda_0=\Lambda/2\Lambda\\
&\text{par un idéal de carré nul}\\
&\text{engendré par }2_\Lambda\ \text{et qui est}\\
&\text{de rang }1\ \text{en tant que }\mathbb F_2\text{-}\\
&\text{module [correspondant donc à}\\
&\text{un \uline{caractère} }\chi:\Lambda_0\to\mathbb F_2].
\end{aligned}\right.
\end{aligned}
\]

Notons d'abord

\uline{Lemme} La \deleted{assertion}\inserted{relation} ``$\forall\,x\in M^*\quad u\wedge x\in\det M^*\Longrightarrow u'\wedge x\in\det M^*$'' équivaut à ``$\forall\,\mathfrak m_i\in\operatorname{Max}(\Lambda)$, \deleted{\unsure{soi}} les images $u_i^\circ$ et $ {u'}_i^\circ$ de $u,u'$ dans le corps résiduel $k_i=\Lambda/\mathfrak m_i$ sont \deleted{proportionnelles}\inserted{colinéaires} [\uline{NB} les conditions $u,u'\in M^*$ \deleted{\unsure{qui}} impliquent que $u_i^\circ,{u'}_i^\circ$ sont $\ne0$] -- c'est \inserted{donc} une relation d'\uline{équivalence} entre $u,u'\in M^*$. Cette relation implique l'existence d'une base $e,f$ de $M^*$, telle que l'on ait $u\wedge e,u\wedge f,u'\wedge e,u'\wedge f\in\det M^*$, i.e.

\[
u=ae+bf,\quad u'=a'e+b'f\qquad\text{avec }a,b,a',b'\in\Lambda^*,
\]

et on peut \unsure{$=$} supposer, quitte à remplacer $e,f$ par $ae,bf$, que $\inserted{a=b=1\quad\text{i.e. }u=e+f}$. Si on suppose maintenant $u\wedge x\in\det M^*\Longrightarrow u'\wedge x\in\pm\{u\wedge x\}$, faisant $x=e$, puis $x=f$, on voit que \deleted{\unsure{$a',b'\in\{\pm a\}$|$a'\in\{\pm a\}$}},

\[
a'\in\{\pm1\},\quad b'\in\{\pm1\},
\]

donc quitte à changer $u'$ de signe, ops

\[
u=e+f,\quad u'=e\pm f
\]
% archive sheet number: 276
